% GENERATED FILE. Edit canonical lessons, then run npm run generate:books.
% canonical-source-hash: fnv1a64:6721eced043e6ddb
% canonical-lessons: GU-C59-fruit, GU-C59-banana, GU-C59-apple, GU-C59-potatoes, GU-C59-tomato

\chapter{Fruit and Vegetables}
\label{ch:gu-fruit-and-vegetables}
\hlchaptermodality{59}

\begin{chapteropening}
\textbf{By the end of this chapter:} \emph{I can name fruit, a banana, an apple, potatoes and a tomato.}

Complete the last lesson of chapter 59: i can name fruit, a banana, an apple, potatoes and a tomato.
\end{chapteropening}

\section[phaḷ]{\gu{ફળ} (phaḷ) — fruit}

\label{lesson:GU-C59-fruit}



\begin{quote}
Before the new one: say the Gujarati for yogurt, then the Gujarati for buttermilk.
\end{quote}

\subsection*{You'll want to know: \gu{ફળ}}
\textbf{\gu{ફળ}} (\emph{phaḷ}) — \textquotedblleft{}fruit\textquotedblright{}. A \textbf{\gu{કેરી}} is one.

\begin{grammarlens}[title={What you've built}]
One more word for fruit and vegetables.
\end{grammarlens}

\subsection*{Your turn}
\begin{itemize}
\raggedright
  \item \emph{Say it:} \emph{phaḷ}
  \item \emph{Say it:} \emph{phaḷ}, once more
  \item \emph{Say it:} say \emph{phaḷ}
  \item \emph{Recall:} say the Gujarati for yogurt, then the Gujarati for buttermilk, then say \emph{phaḷ} again
\end{itemize}

\begin{tcolorbox}[breakable,colback=teal!4,colframe=teal!35!black,title={Before you move on}]
What does \gu{ફળ} mean? (\textquotedblleft{}Fruit\textquotedblright{}.) Say it once more.
\end{tcolorbox}
\section[keḷũ]{\gu{કેળું} (keḷũ) — a banana}

\label{lesson:GU-C59-banana}



\begin{quote}
Before the new one: say the Gujarati for buttermilk, then the Gujarati for fruit.
\end{quote}

\subsection*{You'll want to know: \gu{કેળું}}
\textbf{\gu{કેળું}} (\emph{keḷũ}) — \textquotedblleft{}a banana\textquotedblright{}.

\begin{grammarlens}[title={What you've built}]
One more word for fruit and vegetables.
\end{grammarlens}

\subsection*{Your turn}
\begin{itemize}
\raggedright
  \item \emph{Say it:} \emph{keḷũ}
  \item \emph{Say it:} \emph{keḷũ}, once more
  \item \emph{Say it:} say \emph{keḷũ}
  \item \emph{Recall:} say the Gujarati for buttermilk, then the Gujarati for fruit, then say \emph{keḷũ} again
\end{itemize}

\begin{tcolorbox}[breakable,colback=teal!4,colframe=teal!35!black,title={Before you move on}]
What does \gu{કેળું} mean? (\textquotedblleft{}A banana\textquotedblright{}.) Say it once more.
\end{tcolorbox}
\section[sapharjan]{\gu{સફરજન} (sapharjan) — an apple}

\label{lesson:GU-C59-apple}



\begin{quote}
Before the new one: say the Gujarati for fruit, then the Gujarati for a banana.
\end{quote}

\subsection*{You'll want to know: \gu{સફરજન}}
\textbf{\gu{સફરજન}} (\emph{sapharjan}) — \textquotedblleft{}an apple\textquotedblright{}.

\begin{grammarlens}[title={What you've built}]
One more word for fruit and vegetables.
\end{grammarlens}

\subsection*{Your turn}
\begin{itemize}
\raggedright
  \item \emph{Say it:} \emph{sapharjan}
  \item \emph{Say it:} \emph{sapharjan}, once more
  \item \emph{Say it:} say \emph{sapharjan}
  \item \emph{Recall:} say the Gujarati for fruit, then the Gujarati for a banana, then say \emph{sapharjan} again
\end{itemize}

\begin{tcolorbox}[breakable,colback=teal!4,colframe=teal!35!black,title={Before you move on}]
What does \gu{સફરજન} mean? (\textquotedblleft{}An apple\textquotedblright{}.) Say it once more.
\end{tcolorbox}
\section[baṭākā]{\gu{બટાકા} (baṭākā) — potatoes}

\label{lesson:GU-C59-potatoes}



\begin{quote}
Before the new one: say the Gujarati for a banana, then the Gujarati for an apple.
\end{quote}

\subsection*{You'll want to know: \gu{બટાકા}}
\textbf{\gu{બટાકા}} (\emph{baṭākā}) — \textquotedblleft{}potatoes\textquotedblright{}. A common \textbf{\gu{શાક}}.

\begin{grammarlens}[title={What you've built}]
One more word for fruit and vegetables.
\end{grammarlens}

\subsection*{Your turn}
\begin{itemize}
\raggedright
  \item \emph{Say it:} \emph{baṭākā}
  \item \emph{Say it:} \emph{baṭākā}, once more
  \item \emph{Say it:} say \emph{baṭākā}
  \item \emph{Recall:} say the Gujarati for a banana, then the Gujarati for an apple, then say \emph{baṭākā} again
\end{itemize}

\begin{tcolorbox}[breakable,colback=teal!4,colframe=teal!35!black,title={Before you move on}]
What does \gu{બટાકા} mean? (\textquotedblleft{}Potatoes\textquotedblright{}.) Say it once more.
\end{tcolorbox}
\section[ṭāmeṭũ]{\gu{ટામેટું} (ṭāmeṭũ) — a tomato}

\label{lesson:GU-C59-tomato}



\begin{quote}
Before the new one: say the Gujarati for an apple, then the Gujarati for potatoes.
\end{quote}

\subsection*{You'll want to know: \gu{ટામેટું}}
\textbf{\gu{ટામેટું}} (\emph{ṭāmeṭũ}) — \textquotedblleft{}a tomato\textquotedblright{}.

\begin{grammarlens}[title={What you've built}]
One more word for fruit and vegetables.
\end{grammarlens}

\subsection*{Your turn}
\begin{itemize}
\raggedright
  \item \emph{Say it:} \emph{ṭāmeṭũ}
  \item \emph{Say it:} \emph{ṭāmeṭũ}, once more
  \item \emph{Say it:} say \emph{ṭāmeṭũ}
  \item \emph{Recall:} say the Gujarati for an apple, then the Gujarati for potatoes, then say \emph{ṭāmeṭũ} again
\end{itemize}

\begin{tcolorbox}[breakable,colback=teal!4,colframe=teal!35!black,title={Before you move on}]
What does \gu{ટામેટું} mean? (\textquotedblleft{}A tomato\textquotedblright{}.) Say it once more.
\end{tcolorbox}
